\section{The exact raw density at a minimal flight length}
Let $L_N(x)=S_N\tau_R(x)$ on the collision section, and put
\[
 g=1-2R,\qquad \zeta_R=\operatorname{arcosh}(1+g/R)>0.
\]
Near the orbit alternating normally between centers $0$ and $a$, choose the initial arc coordinate $y$ in the upward tangent direction and the outgoing tangential velocity $p$. Both vanish on the orbit. Use the same upward signed arc coordinate $y_j$ at its successive endpoints.

\begin{lemma}[Reduced action Hessian]\label{lem:hessian}
The second variation of the length with all contact coordinates free is
\begin{equation}\label{eq:quadratic}
 \frac12\sum_{j=0}^{N-1}
 \left\{\frac{y_j^2+y_{j+1}^2}{R}
             +\frac{(y_{j+1}-y_j)^2}{g}\right\}.
\end{equation}
After solving the internal reflection equations, the endpoint Hessian is
\begin{equation}\label{eq:schur}
 S_N=\begin{pmatrix}A_N&-B_N\\-B_N&A_N\end{pmatrix},\quad
 A_N=\frac{\sinh\zeta_R}{g}\coth(N\zeta_R),\quad
 B_N=\frac{\sinh\zeta_R}{g\sinh(N\zeta_R)}.
\end{equation}
In the physical initial coordinates $(y,p)$, $L_N$ has a nondegenerate minimum $Ng$ and Hessian $H_N$ satisfying
\begin{equation}\label{eq:detH}
             \sqrt{\det H_N}=\sinh(N\zeta_R).
\end{equation}
\end{lemma}
\begin{proof}
At successive facing circles the longitudinal endpoint corrections are $y_j^2/(2R)$ and $y_{j+1}^2/(2R)$. Expanding their distance gives \eqref{eq:quadratic}. Its full Hessian is positive definite, since it is a sum of positive square terms including endpoint squares. The interior stationarity equations are
\[
 y_{j+1}-2(1+g/R)y_j+y_{j-1}=0.
\]
Their solution with given endpoints is
\[
 y_j=\frac{\sinh((N-j)\zeta_R)}{\sinh(N\zeta_R)}y_0
       +\frac{\sinh(j\zeta_R)}{\sinh(N\zeta_R)}y_N.
\]
Substitution into the endpoint gradients yields \eqref{eq:schur}. The nonlinear interior equations are uniquely solvable near zero by their positive definite Hessian. Their solution is the physical reflected path by the first variation and the strict incidence at the normal orbit.

The endpoint first variation gives $p=-\partial_{y_0}L_N$ at the linearized level, so $p=-A_Ny_0+B_Ny_N$. The linear map from $(y_0,p)$ to $(y_0,y_N)$ has determinant $1/B_N$. Thus
\[
 \det H_N=\frac{A_N^2-B_N^2}{B_N^2}
                         =\sinh^2(N\zeta_R).
\]
Positivity is preserved under this nonsingular coordinate change, proving the minimum assertion.
\end{proof}

\begin{theorem}[A physical edge coefficient]\label{thm:edge}
Let $\chi$ be a smooth nonnegative section cutoff supported in the regular alternating itinerary, equal to one near its normal initial point. The pushforward of $\chi\nu$ by $L_N$ has, near $Ng$, the form
\begin{equation}\label{eq:edge}
 f_{N,R}(Ng+s)=\one_{\{s\ge0\}}a_{N,R}(s),\qquad
 a_{N,R}(0)=J_{N,R}:=\frac1{2R\sinh(N\zeta_R)}.
\end{equation}
The function $a_{N,R}$ is smooth on a right neighborhood of zero. For each fixed $N$, the construction and smooth bounds may be made uniform on $I$. No assertion of uniform derivative bounds as $N\to\infty$ is made.
\end{theorem}
\begin{proof}
The section density in initial arc and tangential velocity coordinates is $(4\pi R)^{-1}\dd y\dd p$. The Morse lemma applied to Lemma~\ref{lem:hessian} gives $L_N=Ng+|z|^2/2$. The transformed density $w(z)$ is smooth with
\[
 w(0)=\frac1{4\pi R\sqrt{\det H_N}}.
\]
In polar coordinates its pushforward density is
$\one_{\{s\ge0\}}\int_0^{2\pi}w(\sqrt{2s}\,n_\phi)\dd\phi$ near zero. The angular integral of each odd Taylor term vanishes; Taylor's theorem at arbitrary order gives smoothness in $s\ge0$. Its right value is $2\pi w(0)$, which is \eqref{eq:edge}. Compactness in $R$ and the strictly positive determinant give the fixed-$N$ uniform statement.
\end{proof}

\begin{proposition}[Regular critical words and a uniform edge bound]
\label{prop:general-edge}
Fix a regular physical word with $N$ flights between circular obstacles. Every critical point of its total flight length in the two initial section coordinates is normal at both endpoints. There is at most one such point for the fixed sequence of lattice centers. It is a nondegenerate minimum. If $S=\left(\begin{smallmatrix}A&C\\C&D\end{smallmatrix}\right)$ is its reduced endpoint action Hessian, its section-density jump is
\begin{equation}\label{eq:general-edge}
 J=\frac{|C|}{2R\sqrt{AD-C^2}}.
\end{equation}
There is a fixed explicitly computable constant $C_*$ such that, for every regular word and $N\ge2$,
\begin{equation}\label{eq:edge-decay}
                  0<J\le C_*(47/53)^{N-2}.
\end{equation}
No lower incidence margin for the interior collisions is needed for this coefficient bound. This is an individual-word estimate, not a bound for the sum over words.
\end{proposition}
\begin{proof}
Let $q_j(y_j)$ be arc coordinates and $v_j$ the unit velocity of segment $j$. Its second length variation is
\[
 \frac{c_j}{R}y_j^2+\frac{c_{j+1}}{R}y_{j+1}^2+
 \frac{\{v_j^\perp\cdot(t_{j+1}y_{j+1}-t_jy_j)\}^2}{\ell_j},
\]
where $c_j,c_{j+1}>0$ are its outgoing and incoming incidence cosines. At an interior reflected contact the two cosines agree. Summing gives a positive definite tridiagonal matrix with nonzero off-diagonal entries. Its interior Schur complement $S$ is positive definite and $C\ne0$: eliminating the successive interior variables expresses $C$ as the nonzero product of neighboring entries divided by positive pivots.

The endpoint first variation of the reduced action is $-p_0\dd y_0+p_N\dd y_N$. Moreover $\partial_{p_0}y_N=-1/C\ne0$. Thus the derivative of the physical length in $p_0$ vanishes exactly when $p_N=0$, and its other derivative then vanishes exactly when $p_0=0$. At such a critical point the coordinate change gives $\det H=(AD-C^2)/C^2>0$. The same Morse calculation as in Theorem~\ref{thm:edge} proves \eqref{eq:general-edge}.

There cannot be two such regular points for one word. Regard the polygonal length as a convex function of all contact positions in the product of the closed disks. At a normal endpoint its gradient is $-n$, and at an interior reflected contact it is $-2c_j n_j$. These are strictly inward normals. For a different point $x_j$ of the same strictly convex disk, $n_j\cdot(x_j-q_j)<0$. The convex gradient inequality therefore makes every different tuple have strictly larger action. Hence the normal-to-normal critical tuple is the unique global minimizer of that convex problem. Occlusion by a disk not in the word is excluded separately by physical admissibility, not by this minimization argument.

To prove the uniform bound, set $z_j=c_jy_j$ and choose tangent orientations successively so that the off-diagonal entries become negative. At the normal endpoints $c_0=c_N=1$. Write $t_j=1/\ell_j\le t_*=50/3$. The interior matrix has off-diagonals $-t_j$ and diagonals
\[
         d_j=t_{j-1}+t_j+2/(Rc_j)\ge t_{j-1}+t_j+v_*,
             \qquad v_*=200/47.
\]
Its factorization $D_0(I-K)$ has $K\ge0$ with row sums at most
$q=2t_*/(2t_*+v_*)=47/53$. In the Neumann series, an entry connecting interior indices $1$ and $N-1$ vanishes before power $N-2$. Therefore
\[
 |C|\le t_*^2\frac{q^{N-2}}{v_*(1-q)}.
\]
The reduced quadratic form dominates the two endpoint curvature terms, so $S\ge (100/47)I$. Inserting these estimates in \eqref{eq:general-edge}, and using $R\ge9/20$, gives \eqref{eq:edge-decay} with
\[
 C_* = \frac{47}{90}\frac{t_*^2}{v_*(1-q)}.
\]
Arbitrarily small positive interior incidence only increases the diagonal potential, so it does not weaken the estimate.
\end{proof}

The proposition classifies regular critical contributions and bounds each jump uniformly. Singular itinerary boundaries and the number and arrangement of critical values remain to be controlled. Multiplying \eqref{eq:edge-decay} by an uncontrolled word count would not prove a summable global correction.

\begin{corollary}[Raw finite-record absolute continuity]\label{cor:raw-density}
For every $R\in I$ and $n\ge1$, the actual induced record
$(S_n\kappa_R,S_n r_R,S_n\tau_R^Y)$ under $\nu_R^Y$ has a density with respect to counting measure on $\Z^3$ times Lebesgue measure. For fixed total physical count $N$ and displacement $k$, it has the coarea representation, for almost every $t$,
\begin{equation}\label{eq:raw-coarea}
 p_{n,R}(k,N,t)=\frac1{4\pi\nu(Y_R)}
 \sum_{\mathbf d}\int_{D_{n,\mathbf d}\cap\{L_N=t\}}
              \frac{\dd\mathcal H^1}{|\nabla_{\alpha,p}L_N|}.
\end{equation}
Here $D_{n,\mathbf d}$ is the actual first-return event for the ordered center word $\mathbf d$, total count $N$ and displacement $k$. Bounded measurable insertions have the same formula with the insertion in the numerator. Neither a uniform density bound nor a long-time local limit is asserted by this statement.
\end{corollary}
\begin{proof}
Partition the recurrent regular initial states by total collision count and ordered center word. Uniform finite horizon makes the available lattice steps finite; for fixed $N$ there are finitely many words. On every regular branch the physical length is smooth, and Proposition~\ref{prop:general-edge} makes its critical set empty or a singleton. Away from that null set, the coarea formula applied on a countable exhaustion where $|\nabla L_N|$ is bounded below proves absolute continuity and \eqref{eq:raw-coarea}. Restriction to the measurable first-return event does not alter absolute continuity of the initial measure. The singular initial states have zero flux measure, and summing over the countable possible $N$ gives the full law. The same exhaustion with a bounded insertion proves its signed version.
\end{proof}

For the induced system \eqref{eq:Y}, fix $n$ and restrict its joint law to $\kappa=0$, $r=2n$. Every physical flight is at least $g$. Equality of total length with $2ng$ requires every flight to have length $g$, and hence a nearest-neighbor normal alternating orbit. Of the six possible normal starting states only $(0,0)$ lies in $Y_R$. A neighborhood of this point makes exactly $n$ returns of physical length two each. Compactness, or directly the equality cases of the distance bound, excludes other itineraries from a sufficiently small right neighborhood of $2ng$. Consequently this lattice component of the induced law has right density
\begin{equation}\label{eq:inducededge}
                \frac{J_{2n,R}}{\nu(Y_R)}>0
\end{equation}
at its left endpoint, and zero density on its left.

\begin{corollary}[The boundary term must be retained]\label{cor:nonLone}
For the specified inducing set, the raw joint characteristic function $\Phi_{n,R}$ is not in $L^1(\T^3\times\R)$, for any $n\ge1$. In particular a global integrable bound for this raw function cannot be justified merely by making the normal itinerary exponentially rare.
\end{corollary}
\begin{proof}
If it were integrable, taking its $(0,0,2n)$ lattice Fourier coefficient would give an integrable function of the roof frequency. Its inverse would be a continuous density of that component. A continuous representative cannot agree almost everywhere with zero to the left of $2ng$ and with a function having the positive right limit \eqref{eq:inducededge}. This contradicts Theorem~\ref{thm:edge}.
\end{proof}

The central local-limit problem is not decided by this corollary. The endpoint may be far outside the central window, and $J_{2n,R}$ decays exponentially. What changes is the global inversion method, not the raw observable or the intended central theorem.

\begin{proposition}[Exact edge subtraction]\label{prop:subtract}
Choose a smooth roof cutoff supported in the right neighborhood of Theorem~\ref{thm:edge}, equal to one near its endpoint, and denote the resulting subdensity by $f$. With $J=J_{N,R}$, one has the exact identity
\begin{equation}\label{eq:subtract}
 f(Ng+s)=J e^{-s}\one_{\{s\ge0\}}+q(s),\qquad
 \widehat f(b)=e^{ibNg}\left(\frac{J}{1-ib}+\widehat q(b)\right),
\end{equation}
where $q\in L^1(\R)$, its distributional second derivative is a finite signed measure, and
$|\widehat q(b)|\le C_{N,R}(1+b^2)^{-1}$.
\end{proposition}
\begin{proof}
Writing $f(Ng+s)=\one_{s\ge0}\eta(s)a(s)$, the right value of $\eta a-Je^{-s}$ at zero is zero. Thus $q$ is continuous across zero, piecewise smooth, and exponentially decaying at infinity. Its first derivative may jump at zero, so its second distributional derivative has one finite atom and an integrable density. Distributional integration by parts gives $b^2|\widehat q(b)|\le\|D^2q\|_{\TV}$; the $L^1$ norm gives the bound near zero. Direct integration of the exponential proves \eqref{eq:subtract}.
\end{proof}
This is subtraction and exact addition of an explicit density, not convolution with noise or removal of a physical event. It proves an integrable remainder for the isolated edge. Bounds for all remaining branches and edges are still needed for a long-time result.

